\documentclass[a4paper,12pt]{report}  % Déclare un document de type "report" (idéal pour les rapports longs)
% a4paper : Définit le format du papier en A4 (210 × 297 mm)
% 12pt : Définit la taille de police principale à 12 points (standard pour les rapports)

\usepackage[utf8]{inputenc} % Permet d'écrire directement des accents et caractères spéciaux avec l'encodage UTF-8

\usepackage[T1]{fontenc}  % Définit l'encodage des polices pour une meilleure gestion des caractères
% T1 : Active un encodage qui améliore l'affichage des accents (é, à, ç...) 
%      et la césure correcte des mots en français
% Recommandé pour les documents en français

\usepackage{amsmath, amsthm, amssymb}   % amsmath :  Pour des équations et des                                         notations mathématiques avancées.
% amsthm : Pour formater des théorèmes, définitions, propositions, etc.
% amssymb : Pour ajouter des symboles mathématiques supplémentaires

% Définition d'un style pour les définitions
\newtheoremstyle{definitionstyle}  % Nom du style
{10pt}                           % Espace avant
{10pt}                           % Espace après
{\normalfont}                    % Police normale
{}                               % Retrait
{\bfseries}                      % Titre en gras
{.}                              % Ponctuation après le titre
{1em}                            % Espace après le titre (ajout d'un petit                                    espace)
{}                               % Personnalisation du titre

\theoremstyle{definitionstyle}
\newtheorem{definition}{Définition}[chapter] % Numérotation par chapitres

% Définition d'un style pour les exemples
\newtheoremstyle{examplestyle}  % Nom du style
{10pt}                        % Espace avant
{10pt}                        % Espace après
{\normalfont}                 % Police normale
{}                            % Retrait
{\itshape}                    % Titre en italique
{.}                           % Ponctuation après le titre
{1em}                         % Espace après le titre
{}                            % Personnalisation du titre

\theoremstyle{examplestyle}
\newtheorem{example}{Exemple}[chapter] % Numérotation par section

\newtheoremstyle{remarkstyle}  
{10pt}                        % Espace avant
{10pt}                        % Espace après
{\normalfont}                 % Police normale pour le texte
{}                            % Retrait
{\itshape}                    % Titre en italique
{.}                           % Ponctuation après le titre
{1em}                         % Espace après le titre
{}                            % Personnalisation du titre

\theoremstyle{remarkstyle}
\newtheorem{remark}{Remarque}[chapter] % Numérotation par chapitres

% Définition du style pour l'environnement "proposition"
\newtheoremstyle{propositionstyle}  
{10pt}                        % Espace avant
{10pt}                        % Espace après
{\normalfont}                 % Police normale pour le texte
{}                            % Retrait
{\bfseries}                    % Titre en italique
{.}                           % Ponctuation après le titre
{1em}                         % Espace après le titre
{}                            % Pas d'autres spécifications

% Définition de l'environnement proposition
\theoremstyle{propositionstyle}
\newtheorem{proposition}{Proposition}[chapter]  % Numérotation par chapitre

% Définition du style pour l'environnement "théorème"
\newtheoremstyle{theoremstyle}  
{10pt}                        % Espace avant
{10pt}                        % Espace après
{\normalfont}                 % Police normale pour le texte
{}                            % Retrait
{\bfseries}                    % Titre en italique
{.}                           % Ponctuation après le titre
{1em}                         % Espace après le titre
{}                            % Pas d'autres spécifications

% Définition de l'environnement théorème
\theoremstyle{theoremstyle}
\newtheorem{theoreme}{Théorème}[chapter]  % Numérotation par chapitre

% Définition d'un nouveau style de théorème pour les preuves
\newtheoremstyle{proofstyle}  
{10pt}                        % Espace avant
{10pt}                        % Espace après
{\normalfont}                 % Police normale pour le texte
{}                            % Pas de retrait
{\itshape}                    % Titre en italique
{.}                           % Ponctuation après le titre
{1em}                         % Espace après le titre
{}                            % Pas d'autres spécifications

% Appliquer le style à l'environnement de preuve
\theoremstyle{proofstyle}
\newtheorem{preuve}{Preuve}[chapter]  % Numérotation par chapitre

% Définition de l'environnement de preuve avec le carré à la fin
\newenvironment{preuveenv}{\begin{proof}[Preuve]}{\qed\end{proof}}

\usepackage[a4paper, left=2.5cm, right=2.5cm, top=2.5cm, bottom=2.5cm]{geometry} % Définition des marges du document (A4 avec 2.5 cm sur chaque côté)

\usepackage[backend=biber,style=authoryear]{biblatex} % Gestion de la bibliographie
\setlength{\bibitemsep}{1\baselineskip}  % Ajuste l'espacement entre les éléments de la bibliographie

\addbibresource{references.bib} % On inclue le fichier BibTeX references

\usepackage{fancyhdr}     % En-têtes et pieds de page personnalisés

\pagestyle{fancy}         % modification flexible de l'en-tête et du pied de                             page  
\fancyhf{}                % vide l'en-tête et le pied de page
\fancyhead[L]{\leftmark}  % Section ou chapitre courant
\fancyfoot[C]{\thepage}   % Numéro de page au centre

\usepackage{hyperref}
% Ajouter des liens hypertextes dans le document                                 (liens internes, externes, etc.), mettre les liens                             en surbrillance

%%%%%%%%%%%%%%%%%%%%%%%%%%%%%%%%%%%%%%%%%%%%%%%%%%%%
%%%%%%%%%%%%%%%%%%%%%%%%%%%%%%%%%%%%%%%%%%%%%%%%%%%%
%%%%%%%%%%%%%%%%%%%%%%%%%%%%%%%%%%%%%%%%%%%%%%%%%%%%
\usepackage{csquotes}
\usepackage[english,french]{babel} % Langue du document
\usepackage{graphicx}       % Inclusion d'images
\usepackage{float}
\usepackage{color}
\usepackage{comment}        % pour commenter plusieur ligne
\usepackage{enumitem}
\usepackage{titlesec}
\usepackage{calc}
\usepackage{times}
\usepackage{lettrine}
\usepackage{algorithm}
\usepackage{algpseudocode}
\usepackage{mathtools}  % À mettre dans le préambule
\usepackage{listings}
\renewcommand\qedsymbol{$\blacksquare$}

\usepackage{listings}
\usepackage{xcolor}  % pour utiliser les couleurs
\usepackage{ragged2e}

\definecolor{codegreen}{rgb}{0,0.6,0}
\definecolor{codegray}{rgb}{0.5,0.5,0.5}
\definecolor{codepurple}{rgb}{0.58,0,0.82}
\definecolor{backcolour}{rgb}{0.95,0.95,0.92}

\lstdefinestyle{mystyle}{
	backgroundcolor=\color{backcolour},   
	commentstyle=\color{codegreen},
	keywordstyle=\color{magenta},
	numberstyle=\tiny\color{codegray},
	stringstyle=\color{codepurple},
	basicstyle=\footnotesize,
	breakatwhitespace=false,         
	breaklines=true,                 
	captionpos=b,                    
	keepspaces=true,                 
	numbers=left,                    
	numbersep=5pt,                  
	showspaces=false,                
	showstringspaces=false,
	showtabs=false,                  
	tabsize=2
}

\lstset{style=mystyle}

	\lstset{language=Python}
\titleformat
{\chapter} % command
[display] % shape
{\Large} % format
{%
	\rule{\widthof{CHAPITRE}}{0.6pt}\vspace{-0.3cm}\\
	\makebox[0pt][l]{CHAPITRE \Huge\textbf{\thechapter}}\vspace{-0.6cm} \\
	\rule{\widthof{CHAPITRE \Huge\textbf{\thechapter}}}{0.6pt}
} % label
{1ex} % sep
{\huge\textit } 
[
] % after-code

\begin{document}
	
	% En-tête
	\begin{titlepage}
		\begin{center}
			% Logos
			\begin{figure}[t]
				\centering
				\includegraphics[width=0.35\textwidth]{logo/logo_LIRMM.jpg} 
			\end{figure}
			\begin{figure}[t]
				\centering
				\includegraphics[width=0.2\textwidth]{logo/logo_univ_mpt.png} 
			\end{figure}
			
			% Titre
			\vspace{2cm}
			\Large{\textbf{Rapport de stage}}
			
			\vspace{0.5cm}
				
				 \Large{\textbf{Réduction de réseaux - adaptations d'idées provenant du cas polynomiale au cas entier.}} \\
			\vspace{0.5cm}
			\large{\textit{HAI001I : Stage académique}} \\
			
			% Espacement
			\vspace{2cm}
			
			% Auteurs
			\large{
				\textbf{Lucas Noirot} \\
				(\texttt{\href{mailto:lucas.noirot@etu.umontpellier.fr}{lucas.noirot@etu.umontpellier.fr}})
			}
			
			% Espacement
			\vspace{1.5cm}
			
			% Encadrant
			\large{
				\textbf{Encadrant :} \textbf{Romain Lebreton} \\
				(\texttt{\href{mailto:romain.lebreton@lirmm.fr}{romain.lebreton@lirmm.fr}})
			}
			
			% Espacement
			\vspace{2cm}
			
			% Date
			\normalsize{
				\textbf{Date :}  17 février - 26 juin 2025
			}
		\end{center}
	\end{titlepage}
	
	%\chapter*{Remerciements}
	% à rédiger à la fin
	\tableofcontents % Table des matières
	
	\renewcommand{\listalgorithmname}{Table des algorithmes}
	\listofalgorithms % Table des algorithmes
	
	
\chapter*{Notations, complexité et acronymes}
\addcontentsline{toc}{chapter}{Notations, complexité et acronymes}
	
Cette section fait office de lexique et recense l'ensemble des notations et acronymes utilisées tout au long de ce mémoire.
	
\section*{\textit{Notations}}
			
\begin{tabular}{@{}p{3cm}p{12cm}@{}}
	$\mathbb{Z}$ & L'ensemble des entiers relatifs.\\
	$\mathbb{Q}$ & L'ensemble des rationnels.\\
	$\mathbb{R}$ & L'ensemble des réels.\\
	$\mathcal{L}$ & Un réseau, désigné par une lettre majuscule calligraphiée.\\
	$\mathcal{L}($B$)$ & Le réseau engendré par la matrice $B$.\\
	$A^T$ & La transposée de la matrice $A$.\\
	B$^*$ & la base de Gram-Schmidt associée à $B$.\\
	$\mathbb{K}$ & Un corps quelconque.\\
	$\mathbb{F}$ & Un corps fini.\\
	$\mathbb{F}[X]$ & L'anneau des polynômes uni variés à coefficients dans le corps fini $\mathbb{F}$.\\
	$\mathbb{F}_{\leq d}[X]$ & les polynôme sur $\mathbb{F}$ de degré inférieur ou égal $d$.\\
	$\mathbb{F}[X]^{m\times n}$ & les matrices $m\times n$ à coefficients dans $\mathbb{F}[X]$.\\
\end{tabular}
	
\section*{\textit{Complexité}}


\begin{tabular}{@{}p{5cm}p{12cm}@{}}
	Multiplication dans $\mathbb{F}_{\leq d}[X]$ & $M(d) = \mathcal{O}(d \log d \log \log d)$.\\
	Multiplication dans $\mathbb{F}^{n\times n}$ & $MM(n) = \mathcal{O}(n^\omega)$, où $\omega$ est l'exposant de la multiplication matricielle.\\
	Multiplication dans $\mathbb{F}[X]_{\leq d}^{m\times n}$ & $MM(n,d) = \mathcal{O}(\text{MM}(n) M(d)) = \tilde{O}(n^\omega d)$.\\
\end{tabular}

\textbf{Note} : $MM(n,d)$ peut également être obtenu via une approche d'évaluation-interpolation sur une suite géométrique, ce qui permet d'améliorer certaines bornes de complexité.



	\section*{\textit{Acronymes}}
	
	\begin{tabular}{@{}p{3cm}p{10cm}@{}}
		\textbf{LLL} & Lenstra–Lenstra–Lovász\\
		\textbf{BKZ} & Block Korkine-Zolotarev\\
		\textbf{SVP} & Shortest Vector Problem\\
		\textbf{CVP} & Closest Vector Problem
	\end{tabular}

	\chapter*{Guide de Lecture}
	\addcontentsline{toc}{chapter}{Guide de Lecture}
	
		\lettrine{\textbf{D}}{ans} ce mémoire, chaque définition est suivie d'un exemple concret illustrant la notion en question, ainsi que d'un contre-exemple visant à en exposer les subtilités et les exceptions éventuelles. Cette approche permet de mieux comprendre les conditions et les limitations associées à chaque concept. L'objectif est de clarifier les différences entre les situations où une définition est applicable et celles où elle ne l'est pas, afin de renforcer la compréhension approfondie des théorèmes et des constructions présentés.
		
	\chapter*{Introduction}
	
\begin{justify}
	\lettrine{\textbf{L}}{'objectif} de ce stage est d'explorer les réseaux euclidiens et polynomiaux, ainsi que d'analyser les techniques de réduction de réseau. Plus précisément, il s'agit d'adapter les méthodes de réduction développées pour les réseaux polynomiaux au cas des réseaux entiers, afin d'en évaluer l'efficacité et les implications. La réduction de réseau est un outil fondamental en cryptographie, notamment dans les algorithmes de sécurité reposant sur la difficulté de résoudre certains problèmes liés aux réseaux. En particulier, la réduction des réseaux polynomiaux est bien comprise et s'effectue en temps polynomial. L'adaptation de ces techniques aux réseaux entiers pourrait donc avoir un impact significatif sur la sécurité des systèmes cryptographiques, justifiant ainsi une étude approfondie de cette transition et de ses applications pratiques. Par ailleurs, les réseaux euclidiens constituent une approche émergente en cryptographie. Aujourd’hui, la sécurité des systèmes cryptographiques repose principalement sur la difficulté de la factorisation des grands nombres premiers et du calcul de logarithmes discrets. Cependant, ces hypothèses sont menacées par les avancées de l’informatique quantique. À l’inverse, les réseaux euclidiens offrent une alternative prometteuse pour la cryptographie post-quantique, car plusieurs de leurs problèmes fondamentaux semblent résister aux attaques quantiques. Un exemple emblématique est l’algorithme de réduction de bases LLL, qui a permis de briser plusieurs schémas cryptographiques, notamment ceux fondés sur le problème du sac à dos. Pourtant, ce même algorithme joue également un rôle clé dans la conception de nouveaux protocoles sécurisés, en exploitant la complexité intrinsèque des réseaux euclidiens pour garantir la robustesse des schémas cryptographiques.
\end{justify}
	
	
\chapter{Réseaux euclidiens et polynomiaux}

\lettrine{\textbf{L}}{'étude} des réseaux euclidiens en mathématiques trouve ses origines au XVIIIe siècle, lorsque Leonhard Euler a exploré les structures géométriques des points dans l'espace. Toutefois, ce n'est que dans le courant du XXe siècle que le concept de réseaux euclidiens a été intégré à la cryptographie. Dans les années 1990, des chercheurs tels qu'Ajtai, Dwork, Regev et d'autres ont introduit l'idée des réseaux euclidiens comme fondement de problèmes complexes en cryptographie, ouvrant ainsi la voie à de nouvelles constructions cryptographiques.
	
\section{Généralités sur les réseaux}
		
\subsection{Définitions et exemples}
Nous considérons un espace euclidien, c'est-à-dire un espace vectoriel réel de dimension finie muni d'un produit scalaire, noté $\langle f, g \rangle$. Ici, nous utiliserons le produit scalaire usuel, défini par $\langle f, g \rangle = f \cdot g^T$ lequel induit la norme-2, donnée par $\|f\|_2=(f \cdot f^T)^{1/2}$.
	
	
Dans la littérature, on trouve les trois définitions suivantes d’un réseau euclidien qui représentent le même objet.
	
\begin{definition}[\cite{MCA}]
	Soit $n \in \mathbb{N}$ et $f_1, \cdots, f_n \in \mathbb{R}^n$.
		
	Alors $\mathcal{L} = \sum_{1 \leq i \leq n} \mathbb{Z} f_i$ est un réseau euclidien
\end{definition}

\begin{definition}[\cite{Wallet}]
	Un \textbf{réseau euclidien} est un sous-groupe discret de $\mathbb{R}^n$
\end{definition}

\begin{definition}[\cite{clef_unique_2}]
	Un \textbf{réseau euclidien} $(\Lambda, q)$ est un $\mathbb{Z}$-module libre $\Lambda$ de rang fini avec une forme quadratique définie positive $q$ sur $\Lambda \otimes_\mathbb{Z} \mathbb{R}$ 
\end{definition}

\begin{example}
	Les entiers de Gauss forment un réseau dans $\mathbb{C}$ de rang $1$.
\end{example}

\begin{example}
	Un exemple plus exotique, $\mathbb{Z} + \sqrt(2) \mathbb{Z}$ est un réseau de $\mathbb{Z}$.
\end{example}
	
Soit $\mathcal{L}$ un réseau euclidien. Il existe une famille $\mathbb{Z}$-libre maximale $(b_i)_{1 \leq i \leq n}$ dans $\mathcal{L}$ tel que $\mathcal{L} = Zb_1 \oplus \cdots \oplus Zb_n$, qu'on appelle \textbf{base du réseau} $\mathcal{L}$, si on note $B$ la matrice des $(b_i)$ on notera $L(B)$ le réseau de base $B$, donc engendré par les $(b_i)$.
	
L'entier $n$ est commun à toutes les bases de $L$ et on l'appelle rang de $L$. Lorsque $n=m$, on dit que le réseau est de rang plein.
	
\begin{proposition}%Wallet
	Soit $\mathcal{L}$ et $\mathcal{L}'$ deux réseaux de rang $n$ de base $B$ et $B'$.
		
	Alors $\mathcal{L}$ = $\mathcal{L}'$ si et seulement si $\exists U \in M_n(\mathbb{Z})$ tel que $B'=BU$ et $U \in GL_n(\mathbb{Z})$
\end{proposition}

\begin{remark}
	Nous avons la suite d'inclusions suivante : $2\mathbb{Z} \subset \mathbb{Z} \subset \frac{1}{2}\mathbb{Z}$. On observe que $\text{rang}(2\mathbb{Z}) = \text{rang}(\mathbb{Z}) = \text{rang}(\frac{1}{2}\mathbb{Z})$, bien que les ensembles soient distincts, c'est-à-dire $2\mathbb{Z} \neq \mathbb{Z} \neq \frac{1}{2}\mathbb{Z}$. Ce phénomène montre que, contrairement aux espaces vectoriels, pour les réseaux, l'inclusion et avoir le même rang ne suffisent pas à garantir l'égalité des ensembles.
\end{remark}

\begin{remark}
	On remarque également qu'un ensemble de points non alignés dans un réseau ne constitue pas nécessairement une base si son déterminant est différent de $\pm 1$ multiplié par $\|L\|$.
\end{remark}


définition d'un sous réseau

Nous présentons maintenant deux invariants fondamentaux d'un réseau : 
\begin{itemize}
	\item La \textbf{taille du vecteur minimal} du réseau, notée $\lambda_1(\mathcal{L})$, 
	\item Le \textbf{volume du réseau}, aussi appelé la \textbf{taille du réseau} souvent désigné par $|\mathcal{L}|$.
\end{itemize}

\begin{definition}
	La taille d'un réseau $\mathcal{L}(B)$ est $det(B)$ et est noté $\|L\|$. La taille d'un réseau est indépendant de la base choisie.
\end{definition} 

\begin{proposition}
	Soit $\mathcal{L}$ un réseau de $\mathbb{R}^n$. Si $\mathcal{L}'$ est un sous réseau de $\mathcal{L}$ alors $\|\mathcal{L}'\| divise \|\mathcal{L}\|$	
\end{proposition}

En substituant $\mathbb{Z}$ par un anneau de polynômes particulier, on obtient ce qu'on appelle les réseaux polynomiaux. % à détailler car K[X][Y] ne fait pas l'affaire par exemple
	
\begin{definition}[\cite{Wallet}]
	On appelle \textbf{minimum d'un réseau} $\mathcal{L}$ la quantité
		
	$$\lambda_1 = min\{r>0 : |\mathcal{B}(r)\cap\mathcal{L}|>1\} \in \mathbb{R_+}$$
\end{definition}

Plus généralement, pour $k \in \{1, \cdots, n\}$, on pose $\lambda_k(\mathcal{L})$ le plus petit réel $r$ tel qu'il existe $k$ vecteurs $\mathbb{R}$-linéairement indépendants dans $\mathcal{L}$ de norme au plus $r$.

Différentes bases peuvent être associées à un même réseau.
\begin{example}
	Les bases $B = \left(\begin{array}{rr}1 & 0 \\0 & 1\end{array}\right)$ et $B'=\left(\begin{array}{rr}1 & -1 \\-2 & 1\end{array}\right)$ engendrent le même réseau.
		
	Ici $U = B'^{-1} = \left(\begin{array}{rr}-1 & -1 \\-2 & -1\end{array}\right)$ et on voit que $det(U)=-1$. Donc $B$ et $B'$ engendrent le même réseau $\mathcal{L}$.
		
	On a $\|\mathcal{L}\| = 1$ et $\lambda_1( \mathcal{L}) = \lambda_2( \mathcal{L}) = 1$
\end{example}

Existe-t-il une notion de "bonne base" ? Nous verrons qu'une base idéale est celle qui est la plus orthogonale possible. Cette question est en lien avec des problèmes ouverts majeurs, tels que le \textbf{\textit{Shortest Vector Problem}} (SVP) et le \textbf{\textit{closest vector problem}} (CVP).


\begin{theoreme}[Premier théorème de Minkowski]
Pour tout $n \in \mathbb{N}^*$, il existe une constante $C_n > 0$ telle que pour tout réseau $\mathcal{L}$ de $\mathbb{R}^n$, on a :

$$\lambda_1(\mathcal{L}) \leq C_n \|\mathcal{L}\|^{1/n}$$.

En fait, on peut prendre cette constante égale à $C_n=(2/\sqrt{\pi}) \Gamma(n/2+1)^{1/n}$. On appelle constante de Hermite-Minkowski le carré de la constante optimale possible pour cette inégalité, noté $\gamma _n$, en particulier on a $\gamma _n \leq C_n$.
\end{theoreme}

\begin{theoreme}
	Soit $\mathcal{L}$ un réseau de $\mathbb{R}^n$. Soit $S$ une partie convexe symétrique de $\mathbb{R}^n$ telle que $vol(S) > 2^n \|\mathcal{L}\|$, alors il existe $s \in \mathcal{L} \cap S$ non nul. De plus cette inégalité est large si on suppose $S$ compact.
\end{theoreme}

\begin{remark}
	Comme $\|\mathcal{L}\|$ est un invariant du réseau, si on a n'importe quelle base du réseau on peut trouver une approximation de $\lambda_1$.
\end{remark}

\begin{theoreme}[Inégalité d'Hadamard]
Soit $B \in M_n(\mathbb{K})$. On a $\|\det(B)\| \leq \prod_{i=1}^{n} \| b_i \|$, la borne est atteinte si, et seulement si $(b_i)_{1, \cdots, n}$ est une famille orthogonale.	
\end{theoreme}

Une connexion intéressante entre la base de Gram-Schmidt associée à une base d’un réseau et la norme du plus court vecteur de ce réseau.

\begin{proposition}
Soit $\mathcal{L} \subset \mathbb{R}^n$ un réseau de base $f_1, \cdots, f_n$ et soit $f_1^*, \cdots, f_n^*$ sa base de GS associé.

Alors pour tout $x \in \mathcal{L} \backslash \{0\} \|x\| \geq min \{\|f_1^*\|, \cdots, \|f_n^*\|\}$
\end{proposition}


	\subsection{Quelques problèmes algorithmiques liés aux réseaux euclidiens}
	
	Le calcul du plus court vecteur dans un réseau est un problème difficile.
	
	Considérons le problème suivant :
	
	\begin{itemize}
		\item \textbf{Shortest Vector Problem (SVP)} : Étant donnée une base $B$ d’un réseau $L$, trouver un vecteur $v \neq 0$ tel que $\|v\| = \lambda_1(L)$. Ce problème est \textbf{NP-complet} (Ajtai).
	\end{itemize}
	
	On s’intéresse souvent à une version approximative plus accessible :
	
	\begin{itemize}
		\item \textbf{SVP$_\gamma$}, où $\gamma > 0$ : Étant donnée une base $B$ du réseau $L$, trouver un vecteur $v \neq 0$ tel que $\|v\| \leq \gamma \cdot \lambda_1(L)$.
	\end{itemize}
	
	L’état des connaissances actuelles est le suivant :
	
	\begin{itemize}
		\item Pour $\gamma = O(1)$, le problème reste \textbf{NP-complet}.
		\item Pour $\gamma = \text{poly}(n)$, il existe des algorithmes en \textbf{temps exponentiel}.
		\item Pour $\gamma = 2^{O(n)}$, l’algorithme \textbf{LLL} permet de le résoudre en \textbf{temps polynomial}.
	\end{itemize}
	
	Un autre problème important concerne la recherche de vecteurs proches dans un réseau.
	
	\begin{itemize}
		\item \textbf{Closest Vector Problem (CVP)} : Étant donnés une cible $t \in \mathbb{R}^m$ et un réseau $L(B)$, trouver un vecteur $v \in L$ tel que 
		\[
		\|t - v\| = d(t, L) := \min \{ \|t - v\| \mid v \in L \}.
		\]
	\end{itemize}
	
	De même, on peut considérer une version approximative :
	
	\begin{itemize}
		\item \textbf{CVP$_\gamma$}, où $\gamma > 0$ : Trouver un vecteur $v \in L$ tel que 
		\[
		\|t - v\| \leq \gamma \cdot d(t, L).
		\]
	\end{itemize}
	
	Le problème CVP est en général difficile pour un réseau arbitraire. Cependant, pour certaines familles spécifiques de réseaux, comme $\mathbb{Z}^n$, des algorithmes en temps polynomial sont connus. La qualité de la base choisie joue un rôle crucial dans la résolution du problème.
		
	\section{Réduction de réseaux euclidiens et polynomiaux}
	
	
	La réduction de réseaux est un outil fondamental en cryptographie et, plus généralement, en calcul formel. La réduction d'un réseau consiste à modifier une base quelconque de ce réseau en une base presque orthogonale. L'intérêt est de trouver de ``bonnes'' bases afin de résoudre divers problèmes.
	
	Nous avons deux résultats distincts concernant la réduction de réseaux euclidiens et polynomiaux :
	
	\begin{itemize}
		\item La réduction de réseaux sur $\mathbb{F}[X]$ s'effectue en temps polynomial.
		\item La réduction de réseaux sur $\mathbb{Z}$ est $\mathsf{NP}$-difficile.
	\end{itemize}	
	
	\subsection{Réduction de réseaux euclidiens}

	\subsubsection{Généralités}
	
\begin{definition}{\cite{MCA}}
	Soit $f_1, \cdots, f_n \in \mathbb{R}^n$ une base et $f_1^*, \cdots, f_n^*\in\mathbb{R}^n$ sa base de Gram-Schmidt associée.
	
	On dit que $(f_1, \cdots, f_n)$ est \textbf{réduite} si $\|f_i^*\|^2 \leq 2 \| f_{i+1}^*\|^2$ pour $1 \leq i \leq n$ 
\end{definition}
	
\begin{remark}
	Chaque vecteur de la base réduite a une norme au moins égale à la moitié de celle du précédent, garantissant ainsi une décroissance modérée.
\end{remark}	
	
	\begin{comment}
	
	idée : encrypter un truc ave cLLL
	allini def : Une base bi de R n est dite propre si les coeffs de la matrice U apres GS sotn <=1/2
	
	LLL permet de fournir un algo en temps polynomial pour factoriser des polynomes univariés à coefficents entiers. On pourra modifier l'algo pour elf aire factoriser des polynomes bivcariés sur un corps qui a deja la factoriastion sur un pol univarié Q ou cirps finis
	
	parler du fait que le probleme du plsu court vecteur est np dur avec pas d'instance facile par lensa je crois ou arji, on va pouvoir calculer un relativement petit vecteur qui depend de la dimension mais pas dur eseau
	
	
	
	niz d'abeile pas un reseau
	
	intéret a etudier les Q modules ? une sorte d'entre deux entre Z module et R module, 
	
	un resaeu ne peut pas etre dense
	
	En cryptographie il est coutume de considérer des vecteurs en lignes, parler de oplus  de reseau
	
	parler du cout de LLL mca
	
	parler un peu des resaus sans genertaurs genre dual
	
	ca veut dire quoi un polynome modulo m ?
	
	resultant de deux polynomes
	
	parler de O tildechercher un premier lien entre pol et nombre parler de matrice gramiene
	
	\end{comment}

notation pour l'entier le plus proche avec la partie entiere + 1/2	

L'algorithme LLL (proposé par Lenstra et al. en YYYY) repose sur un principe simple : il calcule la meilleure approximation entière de la décomposition de Gram-Schmidt et réduit la base en réorganisant les vecteurs si nécessaire. Cet algorithme est décrit dans \cite{MCA}.	
		
	Le code complet de l'algorithme est présenté dans l'annexe \ref{sec:annexe_code}.
	\begin{algorithm}
		\caption{BasisReduction (LLL) \cite{MCA}}
		\label{algo:LLL_MCA} 
		\ref{algo:LLL_MCA}
			
		\begin{algorithmic}[1]
			\State \textbf{Entrée :} Une base $B=(f_1, \cdots, f_n)$
			\State \textbf{Sortie :} Une base réduite $G=(g_1, \cdots, g_n)$ de $B$
			\For{$i = 1$ \textbf{à} $n$}
				\State $g_i \coloneqq f_i$
			\EndFor
			
			\State $(B^*, U)\coloneqq$ GSO$(B)$
			
			\While{$i \leq n$}
			\For{$j=i-1, i-2$ \textbf{à} $1$}
			\State $g_i:=gi-\lfloor \mu_{ij} \rceil g_j$
			\State Mettre à jour $(B^*, U)\coloneqq$ GSO$(B)$
			\EndFor
			\If{$i>1$ et $\|f_i^*\|^2 > 2 \| f_{i+1}^*\|^2$}
			\State échanger $g_{i-1}$ et $g_i$
			\State Mettre à jour $(B^*, U)\coloneqq$ GSO$(B)$
			\State $i \coloneqq i-1$ 
			\Else 
			\State $i \coloneqq i+1$
			\EndIf
			\EndWhile
			\State Retourner $G=(g_1, \cdots, g_n)$
		\end{algorithmic}
	\end{algorithm}
	
	\subsection{Réduction de réseaux polynomiaux}
	
	
	La réduction des réseaux polynomiaux est un outil essentiel. Par exemple, elle peut être utilisée pour le décodage des codes de Reed-Solomon généralisés. Dans cette partie, nous présenterons les idées et les outils permettant de réduire les réseaux polynomiaux en temps polynomial, en nous appuyant sur les meilleurs exposants connus à ce jour. Lorsque l'on travaille avec des matrices à coefficients dans un corps fini \(\mathbb{F}\), de nombreuses opérations ont des complexités équivalentes comme la multiplication, l'inversion, le calcul du déterminant ou encore la résolution d'un système linéaire. Mais qu'en est-il lorsque les matrices ont leurs coefficients dans \(\mathbb{F}[X]\) ? Dans ce cas, le calcul du déterminant est équivalent à la multiplication de matrices. D'autres opérations, comme l'ordonnancement des bases et la réduction de colonnes, restent également de complexité comparable. En revanche, l'inversion d'une matrice ne l'est plus, en raison de la taille de la sortie. L'objectif de cette partie est donc d'une part de remédier à ces problèmes, mais d'autre part de faire l'état de l'art actuel sur la réduction de réseaux polynomiaux. Cette section sur la réduction des réseaux polynomiaux est fortement inspirée de \cite{CCLebreton}.  
	
	Soit $F \in \mathbb{F}[X]^{m \times n}$. On définit $(F, \sigma):= \{v \in \mathbb{F}^{1 \times m}$ tel que $ vF = 0$ mod $x^\sigma\}$. On appellera sigma la précision
	
	\begin{proposition}
		$(F, \sigma)$ est un $\mathbb{F}[X]$-module de dimension $m$.
	\end{proposition}

	faire la preuve, Il faut montrer que c'est un sous-groupe additif stable pour la multiplication.
	
	\begin{definition}
		Une $(F, \sigma)$ base d'ordre $P$ est une base(au sens des F[X]-modules) $(F, \sigma)$ de degré minimale.
	\end{definition}
	
	c'est quoi degré ? c'est quoi minimale ?
	
	
	\begin{definition}[\cite{clef_unique_0}]
		Soit $\mathbb{K}$ un corps et $M \in \mathbb{K}[x]_{1 \times n}$. On définit le \textbf{degré de ligne} du vecteur ligne M par : 
		$$rdeg(M)=max(deg(m_i))_{i\in\{1, \cdots, n\}}$$
	\end{definition}
	
	\begin{definition}[\cite{clef_unique_0}]
		Soit $\mathbb{K}$ un corps et $M \in \mathbb{K}[x]_{n \times n}$. On définit le \textbf{degré de ligne} de la matrice M par :
		$$rdeg(M)=max(rdeg(ligne_i))_{i\in\{1, \cdots, n\}}$$
	\end{definition}
	
	\begin{example}
		Soit $M =
		\left(\begin{array}{rrrr}
			1 & 0 & 1 & 1 \\
			x & 1 & x + 1 & 0 \\
			1 & x^{3} + x^{2} & x & 0 \\
			x^{2} & 0 & x^{4} + x^{3} & 0
		\end{array}\right) \in \mathbb{F}_2[x]
		$
		alors $row\_degree(M)=\left(\begin{array}{rrrr}
			0 & 1 & 3 & 4
		\end{array}\right)\in \mathbb{Z}^4$
		
	\end{example}
	
	Mais on a un problème, si c = bA, alors rdeg(c) n'est pas lié à rdeg(b)e et rdeg(A) en général
	
	\begin{definition}
		Soit $s \in \mathbb{Z}^n$. On définit le degré de ligne décalé du vecteur ligne M par :
		$$rdeg_s(M)=max(deg(m_i)+s_i\\)$$
	\end{definition}

	degré décalé d'une matrice
	
	Remarque : Onnnote x^s par la matrice diagonale x^s1, ..., xŝn , alors rdeg_s(A) = redeg (A x^s)

	
	exemple
	Si $F =
	\left(\begin{array}{rrrr}
		1 & 0 & 1 & 1 \\
		x & 1 & x + 1 & 0 \\
		1 & x^{3} + x^{2} & x & 0 \\
		x^{2} & 0 & x^{4} + x^{3} & 0
	\end{array}\right) \in \mathbb{F}_2[x]
	$ et $s =(1, 0, 0, 1)$ alors 
	
	rdeg_s F = rdeg (F . x^s) = rdeg $F =
	\left(\begin{array}{rrrr}
		x & 0 & 1 & x \\
		x^2 & 1 & x + 1 & 0 \\
		x & x^{3} + x^{2} & x & 0 \\
		x^{3} & 0 & x^{4} + x^{3} & 0
	\end{array}\right) \in \mathbb{F}_2[x]
 =(1, 2, 3, 4)$
	
	remarque 2 : maintenant tdeg_s(A) = x^-v A x^s
	
	redegsA = v ssi rdeg .. = 0
	redeg <= v ssi <= 0
	
	exemple 
	
	transitivité  des degrés de ligne décalé 14:27
	
	On va définir un ordre sur les row degree
	
	\begin{definition}
		
	\end{definition}
	
	Remarque : ce n'est pas un ordre total
	matrice row reduced def
	
	
	pour parler d'un minimum faudrait un ordre total
	
	
	définiton de matrice réduite par ligne
	
	reformulation de la définition
	on définit une base d'ordre comme une base ligné réduite
	
	propostion : il existe une base réduite par ligne P de (F, sigma))
	
	exemple
	
	remarque : on a l'existance, mais pas l'unicité, si on veut l'unicité, on doit avoir une condition en plus : la forme de Popov
	
	definition du leading coerfrcient
	
	exemple
	
	on revisite le lemme de transitivité
	
	preuve de ca
	
	
	généralisation : 26:39
	Lemme
	Si lcoeff(A) est injective (a gauche), alors A est réduite par lignre
	
	preuve
	
	Comment row reduire une matrice ?
	Forme de Weak Popov :
	31:34
	
	COMMENT rendre une matrice weak popov ?
	
	on a l'algo de mulder storjohann
	
	Une suite decrosisante de row degree doit etre stable
	
	lien avec base de grobner
	
	diapao 35:05
	
	Weak popov => row reduced
	
	Complexité de l'algo
	
	O(m^3 d^2) 
	
	on va faire mieux
	
	
	algo DAC
	 je cite \cite{clef_unique_2} et puis \cite{clef_unique_3} et finalement \cite{clef_unique_4}
	
	\chapter{Adaptation de la réduction de réseaux polynomiaux au cas entier}
	
	\chapter{Réseaux définis par relations plutôt que par générateurs}
	\chapter*{Conclusion et perspectives}
	
	\appendix
	
	\titleformat
	{\chapter} % command
	[display] % shape
	{\Large} % format
	{%
		\rule{\widthof{ANNEXE}}{0.6pt}\vspace{-0.3cm}\\
		\makebox[0pt][l]{ANNEXE \Huge\textbf{\thechapter}}\vspace{-0.6cm} \\
		\rule{\widthof{ANNEXE \Huge\textbf{\thechapter}}}{0.6pt}
	} % label
	{1ex} % sep
	{\huge\textit } 
	[
	] % after-code
	
	
	%rajouter des corps ?
	\chapter{Rappels d'algèbre : Groupes, Anneaux et Modules}

	\section{Définition et propriétés des groupes}
	
	
	subsection Généralités
	
	definition
	
	exemple contre exemple
	
	
	\section{Définition et propriétés des anneaux}
	
	\subsection{Généralités}
	
	\begin{definition}
		Un \textbf{anneau} $(A, +, \cdot)$ est un ensemble : 
		\begin{itemize}
			\item muni d'une loi interne $+$ tel que $(M, +)$ est un groupe abélien.
			\item muni d'une loi interne $A \times A \rightarrow A, (a, b) \mapsto a \cdot b$ qu'on notera $ab$ vérifiant les trois propriétés suivantes :
			
			\begin{enumerate}
				\item $\cdot$ distributif sur $+$ :
					\begin{itemize}
						\item $a(b + c) = ab + ac$
						\item $(b+c)a = ba+ca$
					\end{itemize}
				
				
				\item $\cdot$ est associatif
				
				\item $\cdot$ a un élément neutre
			\end{enumerate}
		\end{itemize}
	\end{definition}

	\begin{definition}
		Un anneau $A$ est \textbf{commutatif} si $ab = ba$ pour tout $a$, $b \in A$
	\end{definition}
	
	\begin{definition}
		Un anneau $A$ est \textbf{intègre} si $ab=0 \Rightarrow a = 0$ ou $b = 0$ pour tout $a$, $b \in A$
	\end{definition}

	\begin{definition}
		Soit $A$ un anneau. Un \textbf{idéal} $I$ de $A$ est un sous-ensemble de $A$ tel que
		\begin{enumerate}
			\item $I$ est un sous-groupe additif de $(A, +)$
			\item pour tout $a \in A$, $x \in I$, on a $ax \in I$ et $xa \in I$
		\end{enumerate}
	\end{definition}

	def anneau euclidien
	
	\begin{definition}
		Un idéal est dit \textbf{principal} s'il est engendré par un élément. Un anneau A est dit \textbf{principal} si tous ses anneaux sont principaux.
	\end{definition} 
	
		

	def Anneau noetherien
	
	prop anneau noetherien
	
	
	def anneau factoriek
	
	\begin{theoreme}
		Tout anneau principal est factoriel.
	\end{theoreme}
	
	
	\begin{theoreme}
		Tout anneau euclidien est principal
	\end{theoreme}
	
	\begin{theoreme}
		Tout anneau principal est noetherien
	\end{theoreme}
	
	
	\begin{theoreme}
		ctout anneau commutatif intègre et principal est foactoriel
	\end{theoreme}
	
	
	

	
	
	proposition des implications
	schéma d'inclusion récap pq pas.
	
	\subsection{Anneaux de polynômes}
	
	Dans ce mémoire, les anneaux de polynômes constituent un type d'anneaux particulièrement utile.
	
	\begin{theoreme}
	Si $A$ est un anneau factoriel, alors $A[X]$ est un anneau factoriel.
	\end{theoreme}

	\begin{remark}
		Il s'en suit que $A[X, Y]$ est factoriel et ainsi de suite.
	\end{remark}

	\begin{theoreme}
		Si $A$ est un anneau Noethérien, alors $A[X]$ est un anneau Noethérien.
	\end{theoreme}
	
	\begin{remark}
		Il s'en suit que $A[X, Y]$ est Noethérien et ainsi de suite.
	\end{remark}

	\begin{theoreme}
		Si $A$ est un anneau factoriel, alors $A[X]$ est un anneau intègre.
	\end{theoreme}
	
	\begin{remark}
		Il s'en suit que $A[X, Y]$ est intègre et ainsi de suite.
	\end{remark}

	Si K corps K X principal je crois
	
	\begin{proposition}
		$K[X]$ est principal.
	\end{proposition}

	\begin{proposition}
		$K[X, Y]$ n'est pas principal.
	\end{proposition}
	
	\section{Définition et propriétés des modules} %pq pas exmeple ?
	
	\subsection{Généralités}
	Cette partie se base sur \cite{ring_modules} et \cite{harari_modules}.
	
	La notion de module est la généralisation naturelle de celle d'espace vectoriel. Dans toute la suite $A$ désigne un anneau commutatif.
	
	\begin{definition}
		Un \textbf{$A$-module} $(M, +, \cdot)$ est un ensemble :
		
		\begin{itemize}
			\item muni d'une loi interne $+$ tel que $(M, +)$ est un groupe abélien.
			\item muni d'une loi externe $A \times M \rightarrow M, (a, m) \mapsto a \cdot m$, qu'on notera $am$ vérifiant les quatre propriétés suivantes :
			\begin{enumerate}
				\item Distributivité : $\alpha(m + m') = \alpha m + \alpha m'$
				\item Distributivité : $(\alpha+\beta)m = \alpha m + \beta m$
				\item Associativité : $(\alpha \beta)m = \alpha(\beta m)$
				\item Neutre : $1m = m$
				
				pour tout $\alpha$, $\beta \in A$ et $m$, $m' \in M$
			\end{enumerate}
		\end{itemize}
	\end{definition}
	
	\begin{remark}
		La définition d'un module peut être vue comme une généralisation de celle d'un espace vectoriel. En effet, un module est un espace vectoriel dans lequel l'hypothèse sur $A$ est affaiblie : au lieu d'être un corps, $A$ est un anneau, ici commutatif, on ne se soucie donc pas de la distinction entre modules à gauche ou à droite.
	\end{remark}

	%plus on affaibli les hypothèses plus l'objet est général et donc plus difficile à étudier
	
	\begin{example}
		Soit $n \in \mathbb{N}$, $\mathbb{Z}/n \mathbb{Z}$ est un $\mathbb{Z}$-module.
		
		Z carré peut etr evu comme un Z module (c'est un groupe abelien pour l'addition, )
	\end{example}

	\begin{definition}
		Soit $M$ un $A$-module. Un \textbf{sous-module} $N$ de $M$ est un sous-groupe de $(M, +)$ stable pour la multiplication externe par tout élément de $A$.
	\end{definition}

	% contre exemple au fait d'etre stable pour la multiplication externe ?
	
	%exemple et contre-exemple de modules
	
	un sous module dans Z n est un sous groupe de Z n avec m < n on a forcément stable multiplcaiton et le citer pour dire qu'un sous groupe d'un reseau est forcement un reseau jsp si c'est uile
	
	\begin{definition}
		Un $A$-module $M$ est dit de \textbf{type fini} s'il existe une partie finie $S$ de $M$ tel que $M$ soit engendré par $S$.
	\end{definition}

	\begin{definition}
		Un $A$-module $M$ est dit \textbf{libre} s'il admet une base i.e. si il admet une famille $(x_i)_{i \in I}$ tel que pour tout $x \in M$, $x$ s'écrit de manière unique comme $\sum_{i \in I} \alpha_i x_i$
	\end{definition}

	\begin{proposition}
		Si $M$ est libre et de type fini, alors il admet une base finie. On dit dans ce cas que $M$ est libre de type fini, et toutes les bases de $M$ ont le même cardinal, qu'on appelle rang de $M$.
	\end{proposition}

	\begin{example}
		$\mathbb{Z}/n \mathbb{Z}$ est un $\mathbb{Z}$-module de type fini mais pas libre pour $n \neq 0$.
	\end{example}
		
	\subsection{Modules sur un anneau principal}
	Dans cette partie, $A$ désigne un anneau principal.
	
	\begin{theoreme}
	Si M est un module libre sur un anneau principal A, tout sous-module de M est libre et de rang inférieur ou égal à celui de M.
	\end{theoreme}


	\begin{example}
		Dans le $\mathbb{Z}$-module libre $\mathbb{Z}^2$, le sous-module des couples $(a, b)$ tels que $a \equiv b$ mod $10$ est libre de base $((1,1), (0, 10))$ donc de rang $2$ (comme le module lui-même).
	\end{example}
	
	
	théorie des modules sur des anneaux principaux : on a la suite d'inclusion
	\chapter{Rappels d'algèbre linéaire}
	
	Dans ce mémoire, les matrices sont considérées comme des matrices dont les éléments sont des vecteurs ligne.
	
	% demander à Romain si je met GS ici ou ailleurs
	
	def base orthogonale
	
	\section{Orthonormalisation de Gram-Schmidt (GSO)}
	
	On se considère dans un espace euclidien $\mathbb{R}^n$. Soit \( B = (b_i)_{1 \leq i \leq n} \) une base de \( \mathbb{R}^n \), on peut construire la \textbf{base de Gram-Schmidt (GS)} associée à cette base, que l'on notera \( B^* = (b^*_i)_{1 \leq i \leq n} \) par le procédé récursif suivant :
	
	
	$$b_1^* := b_1$$
	
	
	$$b_i^* := b_i - \sum_{j=1}^{i-1} \mu_{i,j} b_j^*, \quad \text{avec } \mu_{i,j} = \frac{\langle b_i, b_j^* \rangle}{\| b_j^* \|^2}$$
	
	où \( \mu_{i,j} \) est appelé \textbf{coefficient de Gram-Schmidt} associé au vecteur \( b_i \). \\
	
	Notons \( B^* \) la matrice des \( b_i^* \), et \( U = \left(\begin{array}{rrrr}
		1 & 0 & \cdots & 0\\
		\mu_{2,1} & \ddots & \ddots & \vdots\\
		\vdots & \ddots & \ddots & 0\\
		\mu_{n,1} & \cdots & \mu_{n,n-1} & 1
	\end{array}\right)\). 

Alors, on a la relation : $B = UB^*$
	
	 \begin{remark}
	 	\begin{itemize}
	 		\item  $U$ est triangulaire avec un déterminant égal à $1$.
	 		\item Contrairement à ce qui est souvent pratiqué dans la littérature, nous n'effectuons pas de normalisation, car cela introduirait des racines, ce qui pourrait entraîner un passage des coefficients \( \mathbb{Q} \) à des coefficients \( \mathbb{R} \setminus \mathbb{Q} \). De plus, la normalisation impose un volume égal à 1, ce qui entraîne une perte de l'information volumétrique du réseau associé.
	 		
	 	\end{itemize}
 	\end{remark}
	
	
	\begin{proposition}
		 Soit $B = (b_i)_{1 \leq i \leq n}$ une base de $\mathbb{R}^n$, et $B^* = (b^*_i)_{1 \leq i \leq n}$ sa base de Gram-Schmidt associée.
		 
		 On a $det(B)=det(B^*) = \prod \| b_i^*\|$
	\end{proposition}	

	\begin{proof}
		Nous utilisons le fait que \( U \) est une matrice triangulaire dont les coefficients diagonaux sont égaux à 1, et que \( B^* \) est une base de vecteurs orthogonaux. Ainsi, 
		\[\det(B) = \det(UB^*) = \det(U)\det(B^*) = \det(B^*) = \prod \| b_i^*\|.\]
	\end{proof}
	
	
	%	\item Propriété du plus petit vecteur** : \\lambda_1(\mathcal{L}) \geq \min \| b_i^* \| \)
	
	
	\begin{algorithm}
		\caption{GramSchmidt (GSO)}
		\label{alg:GSO} 
		\ref{alg:GSO}
		
		\begin{algorithmic}[1]
			\State \textbf{Entrée :} Une base $B=(f_1, \cdots, f_n)$
			\State \textbf{Sortie :} La base de Gram-Schmidt associée $B^*=(b^*_1, \cdots, b^*_n)$ et $U$ la matrice des coefficients de GS
			
			\For{$k = 1$ \textbf{à} $n$}
			\State $b_k^* \coloneqq b_k$
			\For{$j = 1$ \textbf{à} $k$}
			\State $U_{kj}\coloneqq \frac{(b_k*b_j^*)}{\|b_j^*\|^2}$
			\State $b_k^* \coloneqq U_{kj} b_j^*$
			\EndFor
			\EndFor
		\end{algorithmic}
	\end{algorithm}
	\begin{theoreme}
		Le coût de l'algorithme \nameref{alg:GSO} est de \( O(n^3) \) opérations arithmétiques dans \( \mathbb{Q} \).
	\end{theoreme}
	
	
	\chapter{Code SageMath}
	Le code source complet est également disponible sur \url{https://github.com/toncompte/tonrepo}.
	\label{sec:annexe_code}
	Voici le code complet de l'algorithme implémenté :
	
	\begin{lstlisting}
	def Gram_Schmidt(B):
	n = B.nrows()
	
	B_star = Matrix(QQ, n, n)
	U = identity_matrix(QQ, n)
	
	B_star[0] = B[0]
	
	for k in range(1, n):
	B_star[k] = B[k]
	for j in range(k):g
	U[k, j] = (B[k] * B_star[j]) / (B_star[j] * B_star[j])
	B_star[k] -= U[k, j]*B_star[j]
	return U, B_star
	
	\end{lstlisting}
	
	
	

	
	\printbibliography
	
\end{document}
